% TEMPLATE-MILC-v3.tex - plantilla maestra UNICA de las guias MILC v3.
%
% La usan las cuatro familias de guias, cada una con su driver:
%   · Grados 6-11          -> scripts/build-guias-g11.py
%   · Bebras               -> scripts/build-guias-bebras.py
%   · Semillero            -> scripts/build-guias-semillero.py
%   · Territorio interior  -> scripts/build-guias-territorio-interior.py
%
% Antes habia tres copias casi identicas (~400 lineas iguales, 6 distintas):
% cada mejora habia que aplicarla tres veces, y en la practica no se hacia
% (el motor de recursos vivia solo en dos de ellas). Lo que varia entre
% familias esta parametrizado en 6 placeholders que cada driver llena
% (se nombran aqui SIN su sintaxis real: el builder reemplaza por texto plano,
% tambien dentro de los comentarios, y un valor multilinea rompe el preambulo):
%
%   HEADER_IZQ        encabezado izquierdo de pagina
%   HEADER_DER        encabezado derecho de pagina
%   PORTADA_ETIQUETA  rotulo sobre el numero grande (GRADO/MOMENTO/MODULO)
%   PORTADA_SERIE     linea de serie de la portada
%   PORTADA_ID        identificador de la guia en la portada
%   PORTADA_CIRCULO   el circulito de grado (°), o vacio si no aplica
%   PDF_TITULO        titulo en texto plano (sin \textbf ni comillas TeX) para
%                     los metadatos del PDF (/Title); lo produce lib_xelatex.texto_pdf
%
% Composicion (auditoria 2026-09-03): las bandas de estacion piden espacio
% (\Needspace*) y prohiben el salto justo despues (\nopagebreak), asi no quedan
% huerfanas al pie; las cajas largas (softbox, drawbox, infoband, puentebox) son
% `breakable`, asi una caja de media pagina ya no arrastra todo a la siguiente
% dejando la anterior al 40 %. Todo texto sobre mostaza o verde va en milcNegro
% (blanco sobre #E5B400 da 1,93:1 y sobre #58A923 2,95:1; ninguno pasa WCAG AA).
% El resto de claves las documenta content/guias/_SCHEMA.md. El script valida
% que no queden placeholders sin reemplazar antes de escribir el archivo final.

% Etiquetado PDF (latex-lab): /Lang, estructura y texto alternativo de las
% figuras, para lectores de pantalla. Debe ir ANTES de \documentclass.
\DocumentMetadata{lang=es-CO,tagging=on}
\documentclass[11pt,letterpaper]{article}
% headheight=24pt: la cabecera derecha lleva dos lineas (identidad de la guia
% y estacion actual). headsep se acorta en la misma medida para que el bloque
% de texto quede exactamente donde estaba con headheight=12pt.
\usepackage[margin=2.0cm,headheight=24pt,headsep=13pt]{geometry}
\usepackage[table]{xcolor}
\usepackage{fontspec}
\usepackage[spanish]{babel}
\usepackage{tikz}
% Librerías TikZ para los diagramas de las guías (bloque `recursos:`):
%  · babel  → babel en español vuelve activos caracteres como «>», y TikZ los usa
%             en sus opciones (>=stealth, ->). Sin esta librería, cualquier
%             diagrama con flechas revienta con «\language@active@arg> has an
%             extra }». Debe ir ANTES de que se dibuje nada.
%  · positioning        → sintaxis «below left=2pt and 3pt».
%  · decorations.pathreplacing → llaves ({brace}) para señalar rangos.
\usetikzlibrary{babel,positioning,decorations.pathreplacing}
\usepackage{graphicx}
% Circuitos: el trabajo de electrónica se hace hoy con micro:bit y sus sensores
% integrados (MakeCode, sin cableado), así que no se necesitan esquemáticos.
% Si algún día entran montajes con protoboard, basta descomentar esta línea:
% los diagramas .tex/.tikz declarados en `recursos:` ya se incrustan solos.
% \usepackage{circuitikz}
\usepackage[most]{tcolorbox}
\usepackage{tabularx}
\usepackage{array}
\usepackage{fancyhdr}
\usepackage{titlesec}
\usepackage[document]{ragged2e}  % bandera a la derecha en todo el cuerpo
\usepackage{enumitem}
\usepackage{microtype}
\usepackage{etoolbox}
\usepackage{needspace}
% hyperref al final del preambulo: metadatos del PDF (titulo, autor, idioma,
% familia), marcadores por estacion (\pdfbookmark) y enlaces sin recuadro.
\usepackage[hidelinks,
  pdftitle={Prompting técnico para escritura — las 5 partes del prompt profesional},
  pdfauthor={PhD. Álvaro Cárdenas Orozco},
  pdfsubject={Tecnología e Informática · Grado 10 · Guía 3 · MILC v3},
  pdflang={es-CO},
  bookmarksopen=true,bookmarksnumbered=false]{hyperref}

% Helvetica Neue no trae la flecha U+2192, asi que cada «→» escrito literalmente
% en el YAML se compilaba a NADA, en silencio (462 flechas en 61 guias: rutas del
% tipo «paleta → tipografia → lockup» salian rotas). Mapeamos el caracter a su
% version matematica, que si tiene glifo. Asi el autor puede seguir escribiendo
% «→» comodamente en el YAML.
\usepackage{newunicodechar}
\newunicodechar{→}{\ensuremath{\rightarrow}}
% Mismo problema con los simbolos de marca y vinetas: el «✓» de «una palomita ✓»
% salia como cuadrito vacio en el PDF de todas las guias que lo usaban.
\usepackage{pifont}
\newunicodechar{✓}{\ding{51}}
\newunicodechar{✗}{\ding{55}}
\newunicodechar{✔}{\ding{52}}
\newunicodechar{•}{\textbullet}
\newunicodechar{≥}{\ensuremath{\geq}}
\newunicodechar{≤}{\ensuremath{\leq}}

\setmainfont{Helvetica Neue}
\setsansfont{Helvetica Neue}
\setlength{\parindent}{0pt}
\setlength{\parskip}{6pt}
% Sin particion de palabras: en 6.º-7.º una palabra cortada por guion al final
% de linea («Comuni-cacion») frena la lectura mucho mas que una linea corta.
\hyphenpenalty=10000
\exhyphenpenalty=10000
\pretolerance=2000
\tolerance=3000
\setlength{\emergencystretch}{3em}
\linespread{1.10}
\renewcommand{\arraystretch}{1.25}
\setlist{leftmargin=5mm,itemsep=1mm,topsep=1mm}
\newcolumntype{Y}{>{\RaggedRight\arraybackslash}X}

\definecolor{milcMagenta}{HTML}{E6007E}
\definecolor{milcVino}{HTML}{5A0038}
\definecolor{milcTurquesa}{HTML}{0093A5}
\definecolor{milcVerde}{HTML}{58A923}
\definecolor{milcMostaza}{HTML}{E5B400}
\definecolor{milcOcre}{HTML}{B86600}
\definecolor{milcNegro}{HTML}{241816}
\definecolor{milcGris}{HTML}{F7F7F4}
\definecolor{milcLinea}{HTML}{E8E2DD}
\definecolor{milcGradePrimary}{HTML}{FF2D87}
\definecolor{milcGradeDark}{HTML}{9F1B56}
\definecolor{milcGradeSoft}{HTML}{FFE0EE}
\definecolor{milcGradeAccent}{HTML}{CFE7FF}
\definecolor{milcGradeText}{HTML}{FFFFFF}
\color{milcNegro}

\pagestyle{fancy}
\fancyhf{}
% Cabecera de dos lineas: identidad de la guia arriba y, a la derecha y en
% gris, la estacion en curso (\markright desde \milcestacion). Antes de la
% primera estacion la segunda linea va vacia; el \strut mantiene la altura
% para que la linea de arriba no baile entre paginas.
\lhead{\small Tecnología e Informática · Grado 10\\[-1pt]{\footnotesize\strut}}
\rhead{\small Guía 3 · MILC v3\\[-1pt]{\footnotesize\strut\color{milcNegro!65}\rightmark}}
\cfoot{\small\thepage}
\renewcommand{\headrulewidth}{0pt}

% Sobre mostaza (#E5B400) y verde (#58A923) el texto va en milcNegro; sobre el
% resto de la paleta, en blanco. Se usa dentro de fonttitle/fontupper, donde un
% \color posterior manda sobre coltitle.
\newcommand{\milctintasobre}[1]{%
  \ifboolexpr{test{\ifstrequal{#1}{milcMostaza}} or test{\ifstrequal{#1}{milcVerde}}}%
    {\color{milcNegro}}{\color{white}}}

\newtcolorbox{titlebox}[1]{
  enhanced,colback=#1,colframe=#1,arc=7mm,boxrule=0pt,
  left=7mm,right=7mm,top=3mm,bottom=3mm,
  before skip=8pt,after skip=8pt,
  fontupper=\sffamily\bfseries\Large\color{white}
}

\newtcolorbox{softbox}[3]{
  enhanced,breakable,pad at break=3mm,
  colback=#2,colframe=#1,borderline west={4pt}{0pt}{#1},
  arc=2mm,boxrule=.4pt,left=4mm,right=4mm,top=3mm,bottom=3mm,
  before skip=6pt,after skip=8pt,title={#3},coltitle=white,
  colbacktitle=#1,fonttitle=\sffamily\bfseries\milctintasobre{#1},fontupper=\RaggedRight,
  attach boxed title to top left={xshift=3mm,yshift=-2mm},
  boxed title style={arc=2mm, boxrule=0pt}
}

\newtcolorbox{drawbox}[2]{
  enhanced,breakable,pad at break=4mm,
  colback=white,colframe=#1,arc=4mm,boxrule=1pt,
  left=5mm,right=5mm,top=4mm,bottom=4mm,before skip=8pt,after skip=8pt,
  title={#2},coltitle=white,colbacktitle=#1,fonttitle=\sffamily\bfseries\milctintasobre{#1},
  fontupper=\RaggedRight,
  attach boxed title to top left={xshift=4mm,yshift=-2mm},
  boxed title style={arc=3mm, boxrule=0pt}
}

\newtcolorbox{infoband}[1]{
  enhanced,breakable,pad at break=2mm,colback=#1!8,colframe=#1!50,arc=2mm,boxrule=.5pt,
  left=4mm,right=4mm,top=2mm,bottom=2mm,before skip=4pt,after skip=4pt,
  fontupper=\small\RaggedRight
}

% Puente narrativo entre secciones (contrato editorial Conéctate).
% Da continuidad: "Acabas de X. Ahora vas a Y porque Z."
% Visualmente: banda izquierda en color del grado, texto en itálica pequeña.
\newtcolorbox{puentebox}{
  enhanced,breakable,pad at break=2mm,colback=milcGradeSoft,colframe=milcGradeDark,
  borderline west={3pt}{0pt}{milcGradeDark},
  arc=0mm,boxrule=0pt,left=5mm,right=4mm,top=2.5mm,bottom=2.5mm,
  before skip=4pt,after skip=8pt,
  fontupper=\itshape\small\color{milcNegro}\RaggedRight
}
% La flecha va en modo matematico a proposito: Helvetica Neue NO tiene el glifo
% U+2192, asi que \textrightarrow se compilaba a NADA («Missing character: There
% is no → in font Helvetica Neue») y el puente salia sin flecha. En modo
% matematico la flecha viene de la fuente matematica y si se dibuja.
\newcommand{\puente}[1]{%
  \begin{puentebox}%
    \textcolor{milcGradeDark}{$\rightarrow$}\hspace{2mm}#1%
  \end{puentebox}%
}

\newcommand{\linead}[1]{\noindent\color{milcLinea}\rule{#1}{0.5pt}\color{milcNegro}}
\newcommand{\respuesta}[1]{\par\vspace{2mm}\foreach \n in {1,...,#1}{\noindent\color{milcLinea}\rule{\linewidth}{0.5pt}\par\vspace{5mm}}\color{milcNegro}}
\newcommand{\checkbox}{\textcolor{milcGradeDark}{\fbox{\phantom{x}}}\hspace{5pt}}

% ─── Listas aireadas para las guías socioemocionales ────────────────────
% Componer las verificaciones como «\checkbox texto\\» deja el texto que salta
% de línea debajo del cuadrito y apelmaza el bloque. Estos entornos alinean la
% continuación y airean los ítems. Los usa Territorio Interior; cualquier
% familia puede hacerlo.
\newlist{checklista}{itemize}{1}
\setlist[checklista]{
  label=\textcolor{milcGradeDark}{\fbox{\phantom{x}}},
  leftmargin=9mm, labelsep=3.2mm, itemsep=2.4mm,
  topsep=2.5mm, parsep=0pt, partopsep=0pt, before=\RaggedRight
}
\newlist{pasoslista}{enumerate}{1}
\setlist[pasoslista]{
  label=\textcolor{milcGradeDark}{\sffamily\bfseries\arabic*.},
  leftmargin=8mm, labelsep=3mm, itemsep=2.8mm,
  topsep=1mm, parsep=0pt, partopsep=0pt, before=\RaggedRight
}

\newcommand{\phasecard}[4]{
\begin{tcolorbox}[enhanced,colback=#3,colframe=#2,arc=3mm,boxrule=.5pt,height=3.05cm,valign=center,halign=center,left=1.5mm,right=1.5mm,top=2mm,bottom=2mm]
{\sffamily\bfseries\fontsize{22}{24}\selectfont\textcolor{#2}{#1}}\par
{\sffamily\bfseries\small #4}
\end{tcolorbox}
}

% ── Piezas del rediseno 2026 (legibilidad 6.º-11.º) ───────────────────
% Banda de estacion: reemplaza el titlebox macizo. Titulo a la izquierda,
% dato practico (tiempo/modalidad) a la derecha, en una sola linea.
% Nunca huerfana: pide 9 lineas libres (o salta de pagina rellenando con
% \vfill) y prohibe el salto justo despues de la banda. Nueve y no seis:
% la caja partible que sigue necesita ~80pt (titulo adosado, relleno y dos
% lineas) para arrancar en la misma pagina; con menos, tcolorbox la manda
% entera a la siguiente y la banda se queda sola. El penalty 10000 va
% en `after`, ANTES del espacio posterior: un `after skip` (aunque sea 0pt)
% es un \vskip tras la caja, y ese glue es un punto de corte legal que un
% \nopagebreak escrito despues de \end{tcolorbox} ya no cierra. Cada banda
% deja marcador PDF y actualiza la cabecera derecha.
\newcounter{milcestacion}
\newcommand{\milcestacion}[2]{%
\Needspace*{9\baselineskip}%
\stepcounter{milcestacion}%
\pdfbookmark[1]{#2}{milcestacion\themilcestacion}%
\markright{#2}%
\begin{tcolorbox}[enhanced,colback=#1,colframe=#1,arc=2mm,boxrule=0pt,
  left=5mm,right=5mm,top=2.6mm,bottom=2.6mm,before skip=11pt,
  after={\par\nopagebreak\vskip7pt}]
{\sffamily\bfseries\fontsize{15}{18}\selectfont\milctintasobre{#1}#2}
\end{tcolorbox}}

% Tarjeta de paso numerada: sustituye las tablas «Pilar / Aplicacion», que
% eran formato de consulta docente, no de estudiante. El numero va en un
% circulo sobre el borde para que la secuencia se lea de un vistazo.
\newcommand{\milcpaso}[3]{%
\Needspace{4\baselineskip}%
\begin{tcolorbox}[enhanced,colback=white,colframe=milcLinea,arc=1.5mm,boxrule=.9pt,
  left=14mm,right=4mm,top=3mm,bottom=3mm,before skip=4pt,after skip=4pt,
  fontupper=\RaggedRight,
  overlay={\node[anchor=north west,circle,fill=#2,text=white,inner sep=0pt,
    minimum size=7.4mm,font=\sffamily\bfseries\small]
    at ([xshift=3.6mm,yshift=-3.2mm]frame.north west) {#1};}]
#3
\end{tcolorbox}}

% Casilla marcable de tamano util con lapiz (la anterior era \fbox{\phantom{x}},
% de alto variable segun el texto que la seguia).
\newcommand{\milcmarca}{\framebox[3.8mm]{\rule{0pt}{3.4mm}}}

% Fila marcable de lista suelta (checks de la estacion 1).
\newcommand{\milccheckrow}[1]{%
\par\vspace{1.8mm}\noindent
\makebox[6.4mm][l]{\raisebox{-.55mm}{\textcolor{milcNegro}{\milcmarca}}}%
\parbox[t]{\dimexpr\linewidth-6.4mm\relax}{\RaggedRight #1}\par}

% Celda marcable dentro de la rubrica (tabla de autoevaluacion). Misma
% sangria francesa, pero pensada para vivir dentro de una columna X.
\newcommand{\milcopcion}[1]{%
\makebox[5.2mm][l]{\raisebox{-.55mm}{\textcolor{milcNegro}{\milcmarca}}}%
\parbox[t]{\dimexpr\linewidth-5.2mm\relax}{\RaggedRight #1}}

% ── Recursos visuales de la guía ──────────────────────────────────────
% Declarados en el bloque `recursos:` del YAML y emitidos por el builder.
% \guiaFigura   → imagen rasterizada o PDF (foto, carta celeste, montaje).
% \guiaDiagrama → fuente TikZ/circuitikz (.tex) incrustada como vector:
%                 es la vía para circuitos y esquemáticos editables.
% [#1] = texto alternativo (alt) · #2 = ruta relativa al .tex · #3 = pie
% El alt llega al PDF etiquetado (/Alt) y lo lee el lector de pantalla.
\newcommand{\guiaFigura}[3][]{%
\begin{center}
\includegraphics[width=0.88\linewidth,height=0.38\textheight,keepaspectratio,alt={#1}]{#2}\\[1.5mm]
{\footnotesize\itshape\textcolor{milcNegro!75}{#3}}
\end{center}
}

\newcommand{\guiaDiagrama}[2]{%
\begin{center}
\input{#1}\\[1.5mm]
{\footnotesize\itshape\textcolor{milcNegro!75}{#2}}
\end{center}
}

\begin{document}
\thispagestyle{empty}

% ── Portada ───────────────────────────────────────────────────────────
% Antes: un rectangulo de color a pagina completa. Se veia bien en pantalla,
% pero en la fotocopiadora del colegio se comia el toner y salia gris sucio.
% Ahora la banda ocupa el tercio superior y el resto va en blanco.
\begin{tikzpicture}[remember picture,overlay]
  \path[fill=milcGradePrimary]
    (current page.north west) rectangle ([yshift=-10.6cm]current page.north east);

  \node[anchor=north west,text=milcGradeText] at ([xshift=2.0cm,yshift=-1.55cm]current page.north west)
    {\sffamily\bfseries\fontsize{12}{14}\selectfont GRADO};
  \node[anchor=north west,text=milcGradeText,opacity=.85] at ([xshift=2.0cm,yshift=-2.35cm]current page.north west)
    {\sffamily\bfseries\fontsize{8.5}{10}\selectfont SERIE GUÍAS · TECNOLOGÍA E INFORMÁTICA};
  \node[anchor=north east,text=milcGradeText,opacity=.85,align=right] at ([xshift=-2.0cm,yshift=-1.55cm]current page.north east)
    {\sffamily\bfseries\fontsize{9}{12}\selectfont I.E. Sor María Juliana\\Cartago, Valle del Cauca};

  \node[anchor=west,text=milcGradeText] (gradonum) at ([xshift=1.85cm,yshift=-7.1cm]current page.north west)
    {\sffamily\bfseries\fontsize{112}{112}\selectfont 10};
  % El circulito de grado (°) solo aplica a las guias de grado («11°»); en
  % Bebras y Semillero el numero grande es un momento/modulo, no un grado.
  % Se posiciona relativo al nodo `gradonum`, no en coordenadas absolutas.
  \draw[milcGradeText,line width=1.6pt]
    ([xshift=2.0mm,yshift=-4.5mm]gradonum.north east) circle (.15cm);

  \node[anchor=west,text=milcGradeText,text width=11.4cm,align=left] at ([xshift=1.5cm,yshift=0.5cm]gradonum.east)
    {
     {\sffamily\bfseries\fontsize{9}{11}\selectfont\MakeUppercase{Período 1 · Escritura abierta con IA}}\\[3mm]
     {\sffamily\bfseries\fontsize{21}{25}\selectfont\setlength{\baselineskip}{27pt}Prompting técnico ---\\[4pt]las 5 partes del\\[4pt]prompt profesional}\\[3.5mm]
     {\sffamily\bfseries\fontsize{15}{18}\selectfont Guía 3-10-TIC}
    };
\end{tikzpicture}

\vspace*{9.4cm}
\vfill

{\sffamily\bfseries\fontsize{8}{10}\selectfont\textcolor{milcGradeDark}{LO QUE TE LLEVAS HOY}}

\begin{tcolorbox}[enhanced,colback=white,colframe=milcNegro,arc=2mm,boxrule=1.6pt,
  left=5mm,right=5mm,top=4mm,bottom=4mm,before skip=3pt,after skip=12pt,
  fontupper=\RaggedRight\fontsize{11}{15.5}\selectfont]
Plantilla personal de prompt con las 5 partes, 3 prompts escritos siguiendo esa plantilla y probados en al menos 2 herramientas, y una bitácora comparativa
\end{tcolorbox}

\begin{tcolorbox}[enhanced,colback=milcGris,colframe=milcGris,arc=2mm,boxrule=0pt,
  left=5mm,right=5mm,top=3.5mm,bottom=3.5mm,after skip=12pt]
{\sffamily\bfseries\fontsize{8}{10}\selectfont\textcolor{milcNegro!55}{ESTA GUÍA ES DE}}\\[3mm]
\begin{tabularx}{\linewidth}{X p{3.2cm} p{3.2cm}}
\linead{\linewidth} & \linead{3.0cm} & \linead{3.0cm}\\[-1mm]
{\fontsize{8.5}{10}\selectfont\textcolor{milcNegro!55}{Nombre completo}} &
{\fontsize{8.5}{10}\selectfont\textcolor{milcNegro!55}{Grupo}} &
{\fontsize{8.5}{10}\selectfont\textcolor{milcNegro!55}{Fecha}}\\
\end{tabularx}
\end{tcolorbox}

\vfill

\noindent\textcolor{milcLinea}{\rule{\linewidth}{1pt}}\\[2mm]
{\sffamily\bfseries\fontsize{9.5}{12}\selectfont Autor: Dr. Álvaro Cárdenas Orozco}\\[1mm]
{\sffamily\fontsize{8.5}{11}\selectfont\textcolor{milcNegro!55}{Metodología MILC · apertura ancestral · pensamiento computacional · triángulo Dussel–estoicismo–Floridi}}

\newpage

\section*{Apertura: Prompting técnico para escritura --- las 5 partes del prompt profesional}

\begin{softbox}{milcGradeDark}{milcGradeSoft}{Saber ancestral}
En Santiago-Manoy, en el Alto Putumayo, las tejedoras ingas hacen \textbf{chumbes}: fajas tejidas con pictogramas geométricos que ellas llaman \emph{labores}. Un trabajo de campo con las tejedoras documentó veinticinco chumbes y cuarenta y seis labores (Aldana Barahona y Sánchez Carballo, 2021). Lo que importa aquí no es el dibujo sino el método. Cada labor se planea en la mente antes de tejerse. La tejedora sabe cómo va a quedar la figura antes de pasar el primer hilo, porque en el telar no hay manera de improvisar a mitad de camino. Y un punto olvidado se nota. El chumbe tampoco es adorno: protege el vientre de la mujer. La \textbf{cara de exclusión} la dice una de las tejedoras, Ruby Rodríguez, en 2018: «estamos en amenaza de que tenemos que ser historia». Hay labores de chumbes antiguos cuyo significado ya nadie conoce.\par\smallskip{\footnotesize\textcolor{milcNegro!70}{\textbf{Fuente.} Aldana Barahona, G. M. y Sánchez Carballo, A. (2021). Tejer con la mente: el chumbe inga del Alto Putumayo colombiano como artefacto cultural y mental. \emph{Estudios Atacameños, 67}, e3521. https://doi.org/10.22199/issn.0718-1043-2021-0007}}
\end{softbox}

\begin{softbox}{milcTurquesa}{milcGris}{Saber contemporáneo}
Un \textbf{prompt} es el texto que le escribes a una IA generativa para que produzca una respuesta. Si el prompt es «escríbeme algo sobre la salud mental», lo que llega es texto genérico que no sirve. Los prompts que funcionan tienen cinco partes. El \textbf{rol}, que dice qué profesional debe asumir la IA. El \textbf{contexto}, con la información de fondo del proyecto. La \textbf{tarea}, con un verbo claro y un objeto preciso. El \textbf{formato}, con la extensión y la estructura esperadas. Y las \textbf{restricciones}, que dicen qué evitar. Un prompt de cien palabras bien armado rinde más que mil palabras de conversación con la herramienta.
\end{softbox}

\begin{softbox}{milcMostaza}{milcGris}{Pregunta-puente}
\textit{¿Qué tiene decidido una tejedora inga antes de pasar el primer hilo? ¿Y qué le falta a quien le escribe a una IA «hazme un capítulo» sin haber decidido nada?}
\end{softbox}

\begin{softbox}{milcVerde}{milcGris}{Saber y saber hacer}
Al terminar podrás: (1) \textbf{identificar} las cinco partes del prompt profesional y qué problema resuelve cada una; (2) \textbf{explicar} por qué un prompt sin alguna de esas partes produce respuestas genéricas; (3) \textbf{aplicar} la plantilla para escribir prompts propios, probarlos en más de una herramienta y documentar cuál funcionó mejor.
\end{softbox}



\small
\begin{tabularx}{\textwidth}{>{\bfseries}p{3.0cm}Y}
\rowcolor{milcGradePrimary!12}Autor & Dr. Álvaro Cárdenas Orozco\\
DBA & Producción editorial mediada por IA con criterio ético y técnico (MEN, T\&I)\\
Referentes & Ley 23 de 1982 · Floridi (ética de la IA generativa) · Dussel (autoría con dignidad)\\
Producto final & Plantilla personal de prompt con las 5 partes, 3 prompts escritos siguiendo esa plantilla y probados en al menos 2 herramientas, y una bitácora comparativa\\
\end{tabularx}
\normalsize

\puente{En la sesión 2 cerraste la ficha editorial y la escaleta. Hoy conviertes ese plano en instrucciones que una IA pueda seguir. Sin las cinco partes, la escaleta que armaste no le llega a la herramienta.}

\milcestacion{milcGradeDark}{Ruta MILC: así vamos a aprender}
\begin{tabularx}{\textwidth}{>{\centering\arraybackslash}X>{\centering\arraybackslash}X>{\centering\arraybackslash}X>{\centering\arraybackslash}X}
\phasecard{1}{milcVerde}{milcGris}{Escucha\\[-1mm]\small Observo, escucho y pregunto.} &
\phasecard{2}{milcTurquesa}{milcGris}{Entender\\[-1mm]\small Sistematizo y ordeno.} &
\phasecard{3}{milcMagenta}{milcGris}{Hacer\\[-1mm]\small Construyo y aplico.} &
\phasecard{4}{milcVino}{milcGris}{Cierre\\[-1mm]\small Evalúo y reflexiono.}
\end{tabularx}


\puente{Antes de teorizar, vas a hacer un experimento con dos prompts sobre lo mismo: uno vago y uno completo. La diferencia entre las dos respuestas es el argumento de toda la sesión.}

\milcestacion{milcVerde}{Estación 1 de 3 · Escucha}

\begin{softbox}{milcVerde}{milcGris}{Escena para escuchar}
\textbf{Actividad 1 · IDENTIFICA --- El experimento de los dos prompts} (15 min · individual). (1) Abre una IA generativa. (2) Escribe el prompt vago: «escríbeme un capítulo para mi libro», y lee la respuesta. (3) Escribe ahora el prompt completo, con el rol, tu tema concreto, tu audiencia, tu género, la extensión que quieres y lo que no quieres. (4) Lee la segunda respuesta. (5) Anota en tres líneas qué información fue la que cambió la calidad.
\end{softbox}

{\sffamily\bfseries\fontsize{8}{10}\selectfont\textcolor{milcNegro!55}{MARCA CUANDO LO HAYAS HECHO}}
\milccheckrow{Escribí los dos prompts sobre lo mismo, el vago y el completo.}
\milccheckrow{Comparé las dos respuestas en lugar de quedarme con la segunda.}
\milccheckrow{Identifiqué qué información concreta fue la que cambió el resultado.}

\begin{infoband}{milcVerde}


\textbf{Lo que va al cuaderno (Actividad 1 · IDENTIFICA).} \textbf{Título:} «Actividad 1 --- El experimento de los dos prompts». \textbf{Formato:} un fragmento corto de cada respuesta y tres líneas sobre la diferencia. \textbf{Extensión:} un tercio de página. \textbf{Sabes que terminaste cuando} puedes nombrar qué dato del segundo prompt hizo la diferencia.
\end{infoband}

\puente{Ya viste la diferencia. Ahora vas a entender la \textbf{anatomía del prompt}: las cinco partes, qué problema resuelve cada una y qué pasa cuando falta alguna.}

\milcestacion{milcTurquesa}{Estación 2 de 3 · Entender}

\begin{softbox}{milcTurquesa}{milcGris}{Lo que vas a entender}
\textbf{Actividad 2 · EXPLICA --- Las cinco partes} (20 min · parejas). (1) Con tu pareja, escriban las cinco partes con el problema que resuelve cada una. (2) Escriban qué pasa cuando falta cada una, con un ejemplo. (3) Armen juntos una plantilla en blanco lista para llenar. (4) Cada uno marca la parte que más se le olvida y dice por qué.
\end{softbox}

\milcpaso{1}{milcTurquesa}{Las cinco partes se descomponen así. El \textbf{rol} activa el vocabulario y el tono de un oficio: «actúa como editor de literatura juvenil». El \textbf{contexto} da el proyecto, la audiencia y el género. La \textbf{tarea} dice qué hacer, con un verbo y un objeto. El \textbf{formato} fija extensión y estructura. Las \textbf{restricciones} dicen qué evitar.}
\milcpaso{2}{milcTurquesa}{El patrón completo es rol, contexto, tarea, formato y restricciones. No hace falta ponerlas en ese orden estricto, pero las cinco tienen que aparecer. Si falta alguna, la herramienta rellena el hueco con lo más probable, que casi siempre es lo más genérico. Invertir dos minutos en escribir bien el prompt ahorra media hora de respuestas inservibles.}
\milcpaso{3}{milcTurquesa}{Todo se juega en una pregunta: ¿la IA tiene la información que necesitaría una persona para hacer este trabajo? Si le pidieras a un editor humano «escribe el capítulo 1» y no le dieras nada más, no podría. La herramienta tampoco. El prompt es exactamente esa información puesta por escrito.}
\milcpaso{4}{milcTurquesa}{Algoritmo: (1) define el rol; (2) escribe el contexto con proyecto, audiencia y género; (3) declara la tarea concreta; (4) especifica el formato esperado; (5) lista las restricciones; (6) envía; (7) lee la respuesta con criterio; (8) ajusta el prompt y vuelve a probar.}

\begin{drawbox}{milcTurquesa}{Anatomía del prompt profesional}
\textbf{Rol:} qué profesional asume la herramienta. \\[1mm] \textbf{Contexto:} proyecto, audiencia, género y escaleta. \\[1mm] \textbf{Tarea:} un verbo claro y un objeto preciso. \\[1mm] \textbf{Formato:} extensión, estructura y elementos. \\[1mm] \textbf{Restricciones:} qué evitar, qué tono no usar. \\[1mm] \textbf{Regla:} si falta una parte, la herramienta rellena con lo más genérico.
\end{drawbox}

\begin{softbox}{milcMostaza}{milcGris}{Errores comunes y su corrección}
(1) Prompt sin rol: la respuesta llega con tono de chatbot genérico. (2) Sin contexto: la herramienta no sabe de qué proyecto ni para quién. (3) Tarea vaga: «algo sobre X» produce algo sobre nada. (4) Sin formato: la respuesta llega con cualquier extensión y estructura. (5) Sin restricciones: aparecen todos los lugares comunes, porque nadie pidió evitarlos.
\end{softbox}

\begin{infoband}{milcTurquesa}


\textbf{Lo que va al cuaderno (Actividad 2 · EXPLICA).} \textbf{Título:} «Actividad 2 --- Las cinco partes». \textbf{Formato:} las cinco partes con el problema que resuelve cada una, la plantilla en blanco y la parte que más se te olvida, marcada. \textbf{Extensión:} media página. \textbf{Sabes que terminaste cuando} podrías construir un prompt profesional sin abrir esta guía.
\end{infoband}

\puente{Tienes el método. Toca aplicarlo. Construyes tu plantilla, escribes tres prompts, los pruebas en al menos dos herramientas y documentas qué funcionó y qué no.}

\milcestacion{milcMagenta}{Estación 3 de 3 · Hacer}

\begin{softbox}{milcMagenta}{milcGris}{Cómo vas a construir tu producto}
\textbf{Actividad 3 · APLICA --- Plantilla, tres prompts y bitácora} (40 min · individual). (1) Escribe tu plantilla personal con las cinco partes y un ejemplo de cada una para tu proyecto. (2) Escribe el prompt del borrador del primer capítulo. (3) Escribe el prompt de tres títulos alternativos. (4) Escribe el prompt de la sinopsis de contraportada. (5) Prueba los tres en al menos dos herramientas distintas. (6) Llena la bitácora con una fila por combinación y anota qué cambiarías de cada prompt.
\end{softbox}

\milcpaso{1}{milcMagenta}{Tu plantilla se construye así: escribe los nombres de las cinco partes, pon debajo de cada una un ejemplo de cómo se ve en tu proyecto y deja el espacio en blanco para llenar en cada uso. Es una guía reutilizable, no un prompt fijo.}
\milcpaso{2}{milcMagenta}{Los tres prompts que vas a probar son distintos a propósito: el borrador de un capítulo, tres títulos alternativos y una sinopsis de contraportada. Uno pide texto largo, otro pide opciones y el tercero pide síntesis. Las cinco partes cambian de contenido en cada caso, pero siguen estando las cinco.}
\milcpaso{3}{milcMagenta}{Probar en más de una herramienta no es capricho. Las respuestas al mismo prompt difieren, y compararlas te enseña más sobre el prompt que sobre las herramientas. Anota cuál preferiste y por qué, con una razón que no sea «sonaba mejor».}
\milcpaso{4}{milcMagenta}{La bitácora tiene una estructura mínima: prompt usado, herramienta, resumen de la respuesta, qué funcionó, qué no funcionó y qué cambiarías del prompt. Una fila por cada combinación de prompt y herramienta.}

\begin{drawbox}{milcMagenta}{Checklist antes de entregar}
\checkbox Plantilla personal con las cinco partes y un ejemplo de cada una. \\[1mm] \checkbox Los tres prompts escritos siguiendo la plantilla. \\[1mm] \checkbox Los tres probados en al menos dos herramientas. \\[1mm] \checkbox Bitácora con una fila por combinación. \\[1mm] \checkbox En cada fila está escrito qué cambiarías del prompt. \\[1mm] \checkbox Borrador del primer capítulo guardado para la sesión 4.
\end{drawbox}

\begin{softbox}{milcMagenta}{milcGris}{Plantilla de trabajo}
\textbf{Plantilla.} \textit{Rol, contexto, tarea, formato y restricciones.} \\[1mm] \textbf{Prompt 1.} \textit{Borrador del primer capítulo.} \\[1mm] \textbf{Prompt 2.} \textit{Tres títulos alternativos.} \\[1mm] \textbf{Prompt 3.} \textit{Sinopsis de contraportada.} \\[1mm] \textbf{Bitácora.} \textit{Una fila por prompt y herramienta, con qué cambiarías.}
\end{softbox}

\begin{infoband}{milcMagenta}


\textbf{Lo que va al cuaderno (Actividad 3 · APLICA).} \textbf{Título:} «Actividad 3 --- Plantilla y tres prompts». \textbf{Formato:} la plantilla, los tres prompts completos, la bitácora comparativa y la decisión final con su razón. \textbf{Extensión:} una página. \textbf{Sabes que terminaste cuando} tienes el borrador del primer capítulo listo para iterarlo en la sesión 4.
\end{infoband}

\puente{Lo que sigue resume \emph{qué} entregas hoy: la plantilla, los tres prompts y la bitácora comparativa. El criterio principal es que la bitácora diga qué cambiarías de cada prompt, no solo cuál te gustó más.}

\milcestacion{milcMostaza}{Producto de la sesión}
\textbf{Plantilla personal, 3 prompts profesionales y bitácora comparativa}.
\begin{softbox}{milcMostaza}{milcGris}{Criterios de calidad}
(1) Plantilla con las cinco partes y un ejemplo propio de cada una. (2) Los tres prompts escritos siguiendo la plantilla. (3) Cada prompt probado en al menos dos herramientas distintas. (4) Bitácora con una fila por combinación de prompt y herramienta. (5) Cada fila dice qué cambiarías del prompt, no solo si gustó. (6) Borrador del primer capítulo guardado como subproducto.
\end{softbox}

\puente{Antes de cerrar, mira el prompting desde las cinco dimensiones humanas. Pedir con precisión es respeto por el trabajo propio. Pedir de cualquier manera y quejarse del resultado es empezar mal el oficio.}

\milcestacion{milcVino}{Evaluación liberadora · cinco dimensiones}

\begin{softbox}{milcVino}{milcGris}{Sentido de la evaluación}
Evaluar no es solo revisar una nota. Es reconocer qué aprendí, qué cambió en mí, qué emociones manejé, cómo me relaciono con la ciudad y la comunidad, y qué le dejaré a quien viene detrás.
\end{softbox}

\textbf{Mi reflexión liberadora en cinco dimensiones.} Completa cada frase iniciada:

\small
\begin{tabularx}{\textwidth}{>{\bfseries\RaggedRight\arraybackslash}p{3.4cm}Y>{\RaggedRight\arraybackslash}p{4.6cm}}
\rowcolor{milcVino}\color{white}Dimensión & \color{white}\bfseries Mi respuesta (frame starter) & \color{white}\bfseries Algo que descubrí\\
1. Personal & Lo que más cambió en mí fue \linead{6.5cm} & \linead{4.2cm}\\
2. Emocional & La emoción más fuerte fue \linead{2.5cm} y la regulé \linead{3cm} & \linead{4.2cm}\\
3. Ciudadana & En democracia esto importa porque \linead{6cm} & \linead{4.2cm}\\
4. Local (Cartago) & En mi barrio sería útil para \linead{6.5cm} & \linead{4.2cm}\\
5. Intergeneracional & A mi abuela/abuelo le diría \linead{6.5cm} & \linead{4.2cm}\\
\end{tabularx}
\normalsize

\begin{infoband}{milcVino}
\textbf{Para la próxima sustentación.} Vuelve a esta tabla en seis meses. Verás que las respuestas cambian — no porque lo aprendido cambie, sino porque tú cambias. La evaluación liberadora es una conversación contigo a través del tiempo.
\end{infoband}


\puente{Y para terminar, tres voces sobre pedir bien. Dussel sobre codificar para que otro decodifique, Epicteto sobre decir solo lo necesario, Floridi sobre saber preguntar.}

\milcestacion{milcGradeDark}{Triángulo de pensamiento · cierre filosófico}

\begin{softbox}{milcVino}{milcGris}{Dussel · lente del nosotros}
\textit{«La información por transmitirse deberá codificarse semántica, sintáctica y fonéticamente, para ser, desde la recepción, decodificada fonética, sintáctica y semánticamente como información recibida.»}

{\footnotesize\itshape\color{milcNegro!70}--- Enrique Dussel · Filosofía de la liberación (1977), §4.2.5.3}

Un prompt es codificación pura. Lo que tienes en la cabeza --- el tema, la audiencia, el tono --- tiene que quedar escrito para que del otro lado se pueda decodificar. Lo que no codificaste no llega. Y la herramienta no pregunta: rellena el hueco con lo más probable.

\textbf{De lo que yo tengo claro sobre mi libro, ¿cuánto quedó realmente escrito en el prompt?}
\end{softbox}
\linead{\linewidth}\\[1mm]
\linead{\linewidth}

\begin{softbox}{milcMostaza}{milcGris}{Estoicismo · Epicteto · Enquiridión, 33 (c. 125 d.C.) · cuidado interior}
\textit{«Guarda el silencio cuanto te fuere posible. Nunca digas sino lo que absolutamente es necesario, y en ello emplea las menos palabras que pudieres.»}

{\footnotesize\itshape\color{milcNegro!70}--- Epicteto · Enquiridión, 33 (c. 125 d.C.)}

No es una invitación a escribir prompts cortos. Es una invitación a que no sobre nada. Las cinco partes son necesarias; lo que sobra son los rodeos, las cortesías y las explicaciones de por qué necesitas el texto. Cien palabras que hacen falta rinden más que mil que llenan.

\textbf{En mi prompt más largo, ¿qué frase podría borrar sin que la respuesta cambie?}
\end{softbox}
\linead{\linewidth}\\[1mm]
\linead{\linewidth}

\begin{softbox}{milcTurquesa}{milcGris}{Floridi · lente de la infoesfera}
\textit{«La partida la ganarán quienes «sepan preguntar y responder» (Platón, Crátilo 390c) y, por tanto, sepan qué datos pueden ser útiles y relevantes, y merecen recogerse y cuidarse. (trad. propia)»}

{\footnotesize\itshape\color{milcNegro!70}--- Luciano Floridi · Big data and their epistemological challenge (2012)}

Producir texto dejó de ser lo difícil. Lo escaso es saber qué pedir. Por eso la sesión no se trata de la herramienta sino de tu pregunta: quien sabe qué necesita obtiene algo aprovechable de cualquier modelo, y quien no lo sabe no lo obtiene de ninguno.

\textbf{Si me quitaran la herramienta que uso, ¿sabría igual qué es lo que necesito pedir?}
\end{softbox}
\linead{\linewidth}\\[1mm]
\linead{\linewidth}

\begin{infoband}{milcNegro}
\textbf{Nota para el docente.} Las tres son citas textuales y llevan su referencia debajo. Puedes remitir a la obra en clase.
\end{infoband}

\milcestacion{milcVino}{Mi autoevaluación · marca una casilla por fila}

{\small
\setlength{\extrarowheight}{2.2pt}
\begin{tabularx}{\textwidth}{>{\RaggedRight\arraybackslash\bfseries}p{3.0cm}YYY}
\rowcolor{milcVino}\color{white}\bfseries Criterio & \color{white}\bfseries Lo logré & \color{white}\bfseries En proceso & \color{white}\bfseries Necesito apoyo\\[.6mm]
Comprensión del tema & \milcopcion{Explico las ideas con mis palabras.} & \milcopcion{Reconozco las ideas, falta precisión.} & \milcopcion{Necesito repasar lo básico.}\\[.8mm]
\rowcolor{milcGris}Pensamiento computacional & \milcopcion{Usé los pilares al armar mi producto.} & \milcopcion{Usé algunos pilares.} & \milcopcion{Necesito releer la estación 2.}\\[.8mm]
Aplicación y producto & \milcopcion{Mi producto cumple los criterios.} & \milcopcion{Está armado, con ajustes pendientes.} & \milcopcion{Aún está en construcción.}\\[.8mm]
\rowcolor{milcGris}Reflexión 5 dimensiones & \milcopcion{Respondí cada una con honestidad.} & \milcopcion{Respondí algunas con detalle.} & \milcopcion{Mis respuestas son muy generales.}\\[.8mm]
Triángulo de pensamiento & \milcopcion{Conecté las 3 voces con mi trabajo.} & \milcopcion{Conecté al menos una.} & \milcopcion{No las relaciono con mi proceso.}\\[.8mm]
\end{tabularx}}

\vspace{3mm}
\textbf{Hoy entendí que:}\\[1mm]
\linead{\linewidth}\\[1mm]
\linead{\linewidth}

\textbf{Mi compromiso para la próxima sesión es:}\\[1mm]
\linead{\linewidth}\\[1mm]
\linead{\linewidth}

\begin{softbox}{milcMostaza}{milcGris}{Cita de despedida · Marco Aurelio}
\textit{``El obstáculo en el camino se vuelve el camino. Lo que estorba al camino se vuelve camino.''}\\[1mm]
Cada error de hoy, bien observado, se convierte en pieza del camino que pisarás mañana.
\end{softbox}

\small
\begin{tabularx}{\textwidth}{>{\bfseries}p{3.6cm}Y}
\rowcolor{milcGradeDark!12}Nombre completo: & \linead{10cm}\\
Fecha: & \linead{4cm} \quad Curso/grupo: \linead{4cm}\\
Mi firma: & \linead{9cm}\\
\end{tabularx}
\normalsize

\begin{infoband}{milcGradeDark}
\textbf{Para la próxima sesión.} Llega con la plantilla, los tres prompts probados, la bitácora y el borrador del primer capítulo. En la sesión 4 vas a iterar ese borrador: del primer texto al capítulo terminado hay varias vueltas, y ninguna la hace sola la herramienta.
\end{infoband}

\end{document}
